\scanpage{67}
\mspage{66}

On trouve ainsi une correspondance biunivoque entre les ens de $\unsure{\mathcal{F}|\mathscr{F}}$-supports et de $\Lambda$-supports
\diagram{\begin{tikzcd}[ampersand replacement=\&, column sep=large, row sep=small]
\Sigma_{\mathcal{F}}^\circ \arrow[r, shift left=1.2ex] \& \Sigma_\Lambda^\circ \arrow[l, shift left=1.2ex] \\
A \arrow[r, mapsto] \& A\cap\Lambda=\operatorname{omb}_\Lambda(A) \\
\operatorname{supp}_{\mathcal{F}}\unsure{S|\sigma} \arrow[d, equal] \& S \arrow[l, mapsto] \\
\{F\in\mathcal{F}\mid\operatorname{omb}_\Lambda(F)\subset S\} \&
\end{tikzcd}}
respectant les relations d'ordre (inclusion) et le passage au complémentaire, que je vais noter $\unsure{\mathbb{C}|\complement}_{\mathcal{F}},\mathbb{C}_\Lambda$.

On a envie, pour ``comprendre les supports'', de prendre $\Lambda$ aussi petit que possible. Le choix $\unsure{L|\mathcal{L}}$ s'impose, mais cela demande que l'axiome $\mathrm{AtL2}$ soit satisfait --- ce n'est pas \unsure{évident}. Ainsi, on trouve (changeant des notations, en laissant tomber les exposants des $\Sigma$ \dots)
\[
\Sigma_{\mathcal{F}}\simeq\Sigma_M\simeq\Sigma_L
\]
% margin note written slanted
\margin{\uline{NB} La notion de support ne \unsure{dépend} \unsure{pas} de $\mathcal{F}$ \unsure{que|ou} \unsure{par} \unsure{l'inclusion|l'ensemble} des \unsure{Anatos|Aratos} \unsure{spéci.} Dans le \unsure{cas} $\mathrm{AtL2}$, il ne \unsure{dépend} \unsure{donc} que de $L$, \unsure{et} de la relation $\parallel_L$.}
(sous-entendu à chaque fois : la notion d'ens. ordonné des supports, définie en termes de $\parallel_{\mathcal{F}}$ resp $\parallel_M$ resp $\parallel_L$).

Un cas particulièrement \uline{agréable} est celui où ``$x\parallel_L y$'' est la relation tautologique : $x\ne y$ (p.ex cas $\unsure{\gamma}$ d'un \unsure{Anatos|Aratos} quasi-ensembliste)
